% A LaTeX template for MSc Thesis submissions to 
% Politecnico di Milano (PoliMi) - School of Industrial and Information Engineering
%
% S. Bonetti, A. Gruttadauria, G. Mescolini, A. Zingaro
% e-mail: template-tesi-ingind@polimi.it
%
% Last Revision: October 2021
%
% Copyright 2021 Politecnico di Milano, Italy. NC-BY

\documentclass{Configuration_Files/PoliMi3i_thesis}

%------------------------------------------------------------------------------
%	REQUIRED PACKAGES AND  CONFIGURATIONS
%------------------------------------------------------------------------------

% CONFIGURATIONS
\usepackage{parskip} % For paragraph layout
\usepackage{setspace} % For using single or double spacing
\usepackage{emptypage} % To insert empty pages
\usepackage{multicol} % To write in multiple columns (executive summary)
\setlength\columnsep{15pt} % Column separation in executive summary
\setlength\parindent{0pt} % Indentation
\raggedbottom  

% PACKAGES FOR TITLES
\usepackage{titlesec}
% \titlespacing{\section}{left spacing}{before spacing}{after spacing}
\titlespacing{\section}{0pt}{3.3ex}{2ex}
\titlespacing{\subsection}{0pt}{3.3ex}{1.65ex}
\titlespacing{\subsubsection}{0pt}{3.3ex}{1ex}
\usepackage{color}

% PACKAGES FOR LANGUAGE AND FONT
\usepackage[english]{babel} % The document is in English  
\usepackage[utf8]{inputenc} % UTF8 encoding
\usepackage[T1]{fontenc} % Font encoding
\usepackage[11pt]{moresize} % Big fonts

% PACKAGES FOR IMAGES
\usepackage{graphicx}
\usepackage{svg}
\usepackage{transparent} % Enables transparent images
\usepackage{eso-pic} % For the background picture on the title page
\usepackage{subfig} % Numbered and caption subfigures using \subfloat.
\usepackage{tikz} % A package for high-quality hand-made figures.
\usetikzlibrary{}
\graphicspath{{./Images/}} % Directory of the images
\usepackage{caption} % Coloured captions
\usepackage{xcolor} % Coloured captions
\usepackage{amsthm,thmtools,xcolor} % Coloured "Theorem"
\usepackage{float}

% STANDARD MATH PACKAGES
\usepackage{amsmath}
\usepackage{amsthm}
\usepackage{amssymb}
\usepackage{amsfonts}
\usepackage{bm}
\usepackage[overload]{empheq} % For braced-style systems of equations.
\usepackage{fix-cm} % To override original LaTeX restrictions on sizes

% PACKAGES FOR TABLES
\usepackage{tabularx}
\usepackage{longtable} % Tables that can span several pages
\usepackage{colortbl}

% PACKAGES FOR ALGORITHMS (PSEUDO-CODE)
\usepackage{algorithm}
\usepackage{algorithmic}

% PACKAGES FOR REFERENCES & BIBLIOGRAPHY
\usepackage[colorlinks=true,linkcolor=black,anchorcolor=black,citecolor=black,filecolor=black,menucolor=black,runcolor=black,urlcolor=black]{hyperref} % Adds clickable links at references
\usepackage{cleveref}
\usepackage[square, numbers, sort&compress]{natbib} % Square brackets, citing references with numbers, citations sorted by appearance in the text and compressed
\bibliographystyle{abbrvnat} % You may use a different style adapted to your field

% OTHER PACKAGES
\usepackage{pdfpages} % To include a pdf file
\usepackage{afterpage}
\usepackage{lipsum} % DUMMY PACKAGE
\usepackage{fancyhdr} % For the headers
\fancyhf{}

% Input of configuration file. Do not change config.tex file unless you really know what you are doing. 
% Define blue color typical of polimi
\definecolor{bluepoli}{cmyk}{0.4,0.1,0,0.4}

% Custom theorem environments
\declaretheoremstyle[
  headfont=\color{bluepoli}\normalfont\bfseries,
  bodyfont=\color{black}\normalfont\itshape,
]{colored}

% Set-up caption colors
\captionsetup[figure]{labelfont={color=bluepoli}} % Set colour of the captions
\captionsetup[table]{labelfont={color=bluepoli}} % Set colour of the captions
\captionsetup[algorithm]{labelfont={color=bluepoli}} % Set colour of the captions

\theoremstyle{colored}
\newtheorem{theorem}{Theorem}[chapter]
\newtheorem{proposition}{Proposition}[chapter]

% Enhances the features of the standard "table" and "tabular" environments.
\newcommand\T{\rule{0pt}{2.6ex}}
\newcommand\B{\rule[-1.2ex]{0pt}{0pt}}

% Pseudo-code algorithm descriptions.
\newcounter{algsubstate}
\renewcommand{\thealgsubstate}{\alph{algsubstate}}
\newenvironment{algsubstates}
  {\setcounter{algsubstate}{0}%
   \renewcommand{\STATE}{%
     \stepcounter{algsubstate}%
     \Statex {\small\thealgsubstate:}\space}}
  {}

% New font size
\newcommand\numfontsize{\@setfontsize\Huge{200}{60}}

% Title format: chapter
\titleformat{\chapter}[hang]{
\fontsize{50}{20}\selectfont\bfseries\filright}{\textcolor{bluepoli} \thechapter\hsp\hspace{2mm}\textcolor{bluepoli}{|   }\hsp}{0pt}{\huge\bfseries \textcolor{bluepoli}
}

% Title format: section
\titleformat{\section}
{\color{bluepoli}\normalfont\Large\bfseries}
{\color{bluepoli}\thesection.}{1em}{}

% Title format: subsection
\titleformat{\subsection}
{\color{bluepoli}\normalfont\large\bfseries}
{\color{bluepoli}\thesubsection.}{1em}{}

% Title format: subsubsection
\titleformat{\subsubsection}
{\color{bluepoli}\normalfont\large\bfseries}
{\color{bluepoli}\thesubsubsection.}{1em}{}

% Shortening for setting no horizontal-spacing
\newcommand{\hsp}{\hspace{0pt}}

\makeatletter
% Renewcommand: cleardoublepage including the background pic
\renewcommand*\cleardoublepage{%
  \clearpage\if@twoside\ifodd\c@page\else
  \null
  \AddToShipoutPicture*{\BackgroundPic}
  \thispagestyle{empty}%
  \newpage
  \if@twocolumn\hbox{}\newpage\fi\fi\fi}
\makeatother

%For correctly numbering algorithms
\numberwithin{algorithm}{chapter}

%----------------------------------------------------------------------------
%	NEW COMMANDS DEFINED
%----------------------------------------------------------------------------

% EXAMPLES OF NEW COMMANDS
\newcommand{\bea}{\begin{eqnarray}} % Shortcut for equation arrays
\newcommand{\eea}{\end{eqnarray}}
\newcommand{\e}[1]{\times 10^{#1}}  % Powers of 10 notation

\newcommand{\usesvg}[3]{
    \IfFileExists{#1}
    {\includesvg[inkscapelatex=false, width=#2\textwidth, height=#3\textheight]{#1}}
    {\includesvg[inkscapelatex=false, width=#2\textwidth, height=#3\textheight]{Images/TODO.svg}} 
}

% arg1 = svg file path, arg2 = width, arg3= height, arg4 = caption
\newcommand{\useSvgWithCaption}[4]{
    \begin{figure}[h]
        \centering
        \usesvg{#1}{#2}{#3}
        \caption{#4}
    \end{figure}
}

\newcommand{\useImg}[3]{
    \IfFileExists{#1}
    {\includegraphics[width=#2\textwidth, height=#3\textheight]{#1}}
    {\includesvg[inkscapelatex=false, width=#2\textwidth, height=#3\textheight]{Images/TODO.svg}}
}

% arg1 = label [optional], arg2 = img file path, arg3 = width, arg4 = height, arg5 = caption
\newcommand{\useImgWithCaption}[5][]{%
    \begin{figure}[h]
        \centering
        \useImg{#2}{#3}{#4}%
        \caption{#5}%
        \ifthenelse{\equal{#1}{}}{}{%
            \label{#1}%
        }%
    \end{figure}
}

%----------------------------------------------------------------------------
%	ADD YOUR PACKAGES (be careful of package interaction)
%----------------------------------------------------------------------------
\usepackage{tikz}
\usetikzlibrary{positioning}
%----------------------------------------------------------------------------
%	ADD YOUR DEFINITIONS AND COMMANDS (be careful of existing commands)
%----------------------------------------------------------------------------

%----------------------------------------------------------------------------
%	BEGIN OF YOUR DOCUMENT
%----------------------------------------------------------------------------

\begin{document}

\fancypagestyle{plain}{%
\fancyhf{} % Clear all header and footer fields
\fancyhead[RO,RE]{\thepage} %RO=right odd, RE=right even
\renewcommand{\headrulewidth}{0pt}
\renewcommand{\footrulewidth}{0pt}}

%----------------------------------------------------------------------------
%	TITLE PAGE
%----------------------------------------------------------------------------

\pagestyle{empty} % No page numbers
\frontmatter % Use roman page numbering style (i, ii, iii, iv...) for the preamble pages

\puttitle{
	title=Self-Supervised 3D Human Pose Estimation in Sports Aerial Videos, % Title of the thesis
	name=Angelo Attivissimo, % Author Name and Surname
	course=Computer Science and Engineering - \\Ingegneria Informatica, % Study Programme (in Italian)
	ID  = 233484,  % Student ID number (numero di matricola)
	advisor= Prof. Giacomo Boracchi, % Supervisor name
	coadvisor={Diego Martin}, % Co-Supervisor name, remove this line if there is none
	academicyear={2024-25},  % Academic Year
} % These info will be put into your Title page 

%----------------------------------------------------------------------------
%	PREAMBLE PAGES: ABSTRACT (inglese e italiano), EXECUTIVE SUMMARY
%----------------------------------------------------------------------------
\startpreamble
\setcounter{page}{1} % Set page counter to 1

% ABSTRACT IN ENGLISH
\chapter*{Abstract} 
3D Human Pose Estimation (3D HPE) aims to recover the three-dimensional configuration of the human body from visual data and has achieved remarkable progress in recent years thanks to deep learning and large-scale datasets. In sports, 
accurate 3D pose reconstruction can support performance analysis, injury prevention, and tactical evaluation. However, current state-of-the-art methods are mainly designed for controlled environments or ground-level cameras, and are 
not tailored to aerial videos.

Drone platforms offer a flexible and scalable solution for outdoor sports scenarios, where the monitored area can be very large and the installation of multiple calibrated fixed cameras is impractical or economically unfeasible. 
Unmanned Aerial Vehicles (UAVs) can dynamically adjust viewpoint, height, and coverage, enabling data acquisition in complex environments. Nevertheless, aerial footage introduces additional challenges, including viewpoint variability, motion blur, scale changes, 
and frequent occlusions.

In this thesis, I propose Brancher, a deep learning framework built upon a state-of-the-art backbone for single-person, single-view 3D HPE from sports aerial videos. To the best of our knowledge, this is among the first works to explicitly address 
3D human pose reconstruction from monocular drone footage using a fully deep learning approach. The model follows a self-supervised training strategy and extends standard 3D regression by jointly modeling spatio-temporal joint 
dynamics, limb rotations, and prediction uncertainty. By explicitly accounting for rotational alignment and heteroscedastic uncertainty, Brancher improves robustness to noisy pseudo-labels and limited visibility, enabling stable and 
temporally coherent 3D motion reconstruction from challenging aerial sports footage. These results highlight the potential of uncertainty-aware self-supervision for advancing 3D HPE in real-world UAV applications.
\\
\\
\textbf{Keywords:} Deep Learning, Computer Vision, 3D Human Pose Estimation, UAV Footage, Self-Supervised Learning % Keywords

% ABSTRACT IN ITALIAN
\chapter*{Abstract in lingua italiana}
La stima della posa umana 3D (3D HPE) mira a ricostruire la configurazione tridimensionale del corpo umano a partire da dati visivi e ha compiuto progressi significativi negli 
ultimi anni grazie al deep learning e alla disponibilità di dataset su larga scala. In ambito sportivo, una ricostruzione della posa 3D può supportare l'analisi delle prestazioni, la prevenzione 
degli infortuni e la valutazione tattica. Tuttavia, gli attuali metodi allo stato dell'arte sono progettati principalmente per ambienti controllati o telecamere poste ad altezza uomo, e non risultano 
ottimizzati per le riprese aeree.

Le piattaforme basate su droni offrono una soluzione flessibile e scalabile per scenari sportivi all'aperto, dove l'area monitorata può essere molto vasta e l'installazione di sistemi multicamera fissi e 
calibrati risulta spesso logisticamente o economicamente proibitiva. Gli aeromobili a pilotaggio remoto (UAV) possono regolare dinamicamente il punto di vista, l'altezza e la copertura, consentendo 
l'acquisizione di dati in ambienti complessi. Ciononostante, i filmati aerei introducono ulteriori sfide, tra cui la variabilità del punto di vista, il motion blur, i cambiamenti di scala e le frequenti 
occlusioni.

In questa tesi, si propone Brancher, un framework di deep learning basato su una backbone allo stato dell'arte per la 3D HPE monoculare di un singolo soggetto da video sportivi aerei. Questo lavoro è tra i primi ad affrontare esplicitamente 
la ricostruzione della posa umana 3D da riprese monoculari di droni attraverso un approccio basato sul deep learning. 
Il modello adotta una strategia di addestramento auto-supervisionata ed estende la regressione 3D standard modellando congiuntamente la dinamica spazio-temporale delle articolazioni, le rotazioni degli 
arti e l'incertezza della predizione. Tenendo esplicitamente conto dell'allineamento rotazionale e dell'incertezza eteroschedastica, Brancher migliora la robustezza rispetto a pseudo-label rumorose e alla 
visibilità limitata, consentendo una ricostruzione del movimento 3D stabile e temporalmente coerente da riprese aeree complesse. Questi risultati evidenziano il potenziale dell'auto-supervisione consapevole 
dell'incertezza per l'avanzamento della 3D HPE in applicazioni reali con sistemi UAV.
\\
\\
\textbf{Parole chiave:} Deep Learning, Visione Artificiale, Stima della posa umana 3D, Riprese da UAV, Apprendimento auto-supervisionato % Keywords (italian)

%----------------------------------------------------------------------------
%	LIST OF CONTENTS/FIGURES/TABLES/SYMBOLS
%----------------------------------------------------------------------------

% TABLE OF CONTENTS
\thispagestyle{empty}
\tableofcontents % Table of contents 
\thispagestyle{empty}
\cleardoublepage

%-------------------------------------------------------------------------
%	THESIS MAIN TEXT
%-------------------------------------------------------------------------
% In the main text of your thesis you can write the chapters in two different ways:
%
%(1) As presented in this template you can write:
%    \chapter{Title of the chapter}
%    *body of the chapter*
%
%(2) You can write your chapter in a separated .tex file and then include it in the main file with the following command:
%    \chapter{Title of the chapter}
%    \input{chapter_file.tex}
%
% Especially for long thesis, we recommend you the second option.

\addtocontents{toc}{\vspace{2em}} % Add a gap in the Contents, for aesthetics
\mainmatter % Begin numeric (1,2,3...) page numbering

% --------------------------------------------------------------------------
% NUMBERED CHAPTERS % Regular chapters following
% --------------------------------------------------------------------------

\chapter{Introduction}

\section{Human Pose Estimation}

Human Pose Estimation (HPE) is a fundamental task in \textbf{computer vision} that aims to localize human body joints—such as the head, shoulders, elbows, wrists, hips, knees, and ankles—within images or video sequences. 
These joints are subsequently connected according to a predefined kinematic structure to form a skeletal representation of the human body. 
By providing a compact and semantically meaningful description of human posture and motion, HPE enables machines to interpret and reason about human actions, movements, and behaviors in visual data.

Depending on the target representation, human pose estimation can be formulated either as a \textbf{two-dimensional (2D)} or a \textbf{three-dimensional (3D)} problem.
In 2D human pose estimation, joint locations are predicted in the image plane, typically using mathematical methods that operate directly on RGB images or videos. 
While 2D pose estimation has achieved high accuracy under controlled conditions, it inherently lacks depth information and is therefore insufficient for capturing the full spatial structure of human motion.
In contrast, 3D human pose estimation aims to reconstruct joint positions in three-dimensional space, providing a more complete and physically meaningful representation of the human body.
However, recovering 3D pose from visual data—especially from monocular images—introduces significant challenges due to depth ambiguity and the loss of geometric information during the image projection process.

Historically, accurate 3D human pose acquisition has been achieved through \textbf{marker-based} motion capture systems, which rely on reflective markers attached to the subject's body and multiple calibrated cameras to triangulate 
joint positions.
Although such systems provide highly precise measurements, they require specialized hardware, controlled environments, and intrusive setups, limiting their applicability to laboratory conditions.
More generally, motion capture can also be achieved through \textbf{sensor-based} approaches, which rely on wearable devices and instrumented environments to track body motion with high accuracy.
While these systems are effective in controlled settings, they remain dependent on dedicated infrastructure and do not fall within the scope of vision-based human pose estimation.\\
To overcome the limitations of marker-based and sensor-based systems, \textbf{markerless} pose estimation approaches have been developed, leveraging computer vision and learning-based methods to estimate human pose directly from images or videos without the need for physical markers or wearable devices.
Despite their increased flexibility, markerless methods—particularly in 3D—remain sensitive to occlusions, viewpoint changes, and the availability of annotated training data.

Human pose estimation plays a crucial role in a wide range of application domains, including healthcare and rehabilitation, surveillance and safety monitoring, human-computer interaction, animation, and sports performance analysis.
Among these, sports analysis represents one of the most challenging scenarios, as it involves fast and complex motions, self-occlusions, and highly dynamic viewpoints.
These characteristics make sports an ideal testbed for advancing robust and generalizable human pose estimation techniques, particularly in unconstrained real-world environments.

\useImgWithCaption[]{Images/marker_based_vs_marker_less.png}{0.5}{0.5}{An example of a marker-based system vs a marker-less system from \cite{zago20203d}}

\section{State of the Art and Limitations of Current Approaches}

In recent years, human pose estimation has witnessed substantial progress, largely driven by the rapid evolution of deep learning techniques. 
The growing relevance of this research area is reflected in the extensive body of literature, with millions of publications addressing human pose estimation and closely related topics.
This remarkable expansion highlights both the scientific interest and the practical importance of accurately modeling human posture and motion from visual data.

Modern deep learning-based approaches have significantly improved the performance of human pose estimation systems, particularly in controlled environments.
State-of-the-art architectures, including convolutional and transformer-based models, are capable of producing highly accurate and stable pose predictions when operating under favorable conditions.
Such conditions typically include limited occlusions, well-defined and relatively static camera viewpoints, consistent lighting, and the availability of large-scale annotated datasets.
Under these assumptions, both 2D and 3D pose estimation methods have achieved impressive results, often approaching human-level accuracy in benchmark datasets.

Despite these advances, the strong performance of current methods is often tightly coupled to the constraints under which they are developed and evaluated.
Many state-of-the-art models rely on assumptions that are difficult to satisfy in real-world scenarios, such as fixed or frontal viewpoints, minimal camera motion, and the presence of training data closely matching the target domain.
As a result, the generalization capabilities of these models remain limited when deployed in unconstrained environments.

Several challenges remain open and continue to hinder the robustness of human pose estimation systems.
Occlusions caused by self-intersections or external objects can lead to missing or inaccurate joint predictions.
Viewpoint variability introduces significant appearance changes, making it difficult for models trained on canonical perspectives to generalize effectively.
In monocular settings, depth ambiguity represents a fundamental limitation, as multiple three-dimensional poses may correspond to the same two-dimensional projection.
Furthermore, the scarcity of accurately annotated 3D pose data—especially in outdoor and in-the-wild conditions—poses a major obstacle to supervised learning approaches.

As the field of human pose estimation has reached a high level of maturity, further progress increasingly depends on addressing specific and challenging application scenarios where existing methods exhibit clear limitations.
Rather than pursuing incremental improvements on well-established benchmarks, contemporary research efforts focus on identifying underexplored problem settings that better reflect real-world conditions.
Within this context, constrained yet practically relevant scenarios—such as monocular 3D pose estimation in dynamic outdoor environments—represent a meaningful direction for advancing the applicability and robustness of human pose 
estimation systems.

\section{Human Pose Estimation in Outdoor Sports}

Sports performance analysis represents one of the most demanding and impactful application domains for human pose estimation.
Accurate modeling of an athlete’s body motion enables detailed biomechanical analysis, quantitative assessment of technical execution, performance optimization, and injury prevention.
For coaches, athletes, and sports scientists, access to reliable pose information provides valuable insights into movement efficiency, coordination patterns, and load distribution during complex athletic actions.

Traditionally, high-precision motion analysis in sports has been achieved through marker-based motion capture systems.
While these systems offer accurate 3D measurements, they require reflective markers to be attached to the athlete’s body, specialized equipment, and controlled acquisition environments.
In many sports scenarios, the use of physical markers is impractical or undesirable, as it may interfere with natural movement, restrict the athlete’s performance, or be incompatible with competitive settings.
For this reason, markerless pose estimation approaches are particularly attractive for sports applications, as they allow motion analysis to be conducted without altering the athlete’s behavior.

In indoor sports environments, such as volleyball, basketball, or tennis played in controlled arenas, effective markerless solutions have been developed by leveraging multi-camera systems.
By using multiple synchronized and calibrated cameras, accurate 3D pose reconstruction can be achieved through geometric triangulation.
The fixed camera setup, limited capture volume, and controlled lighting conditions make these environments well-suited for multi-view pose estimation techniques.
As a result, reliable 3D motion analysis has become feasible in several indoor sports contexts.

However, extending these approaches to \textbf{outdoor} sports introduces substantial challenges.
Outdoor sports often take place over large spatial areas, making the installation and calibration of dense networks of static cameras logistically complex and economically costly.
Environmental factors such as varying illumination, weather conditions, and background clutter further complicate visual data acquisition.
Moreover, athletes may move freely across wide regions, frequently leaving the field of view of fixed cameras.
As a consequence, traditional multi-camera motion capture systems struggle to scale beyond controlled indoor environments and fail to capture natural athletic movements in real-world outdoor settings.

These limitations highlight the need for flexible and non-intrusive pose estimation solutions capable of operating in unconstrained outdoor sports scenarios, where controlled camera setups are not feasible.

\section{Drones as a Data Acquisition Platform}

The limitations of fixed camera systems in outdoor sports scenarios motivate the adoption of mobile sensing platforms capable of flexibly following athletes across large and unconstrained environments.
In this context, unmanned aerial vehicles (UAVs), commonly referred to as drones, have emerged as a promising solution for visual data acquisition in sports analysis.
By providing a mobile and adaptive viewpoint, drones enable the capture of athletic performance in settings where static camera installations are impractical or infeasible.

Over the past decade, drone technology has undergone substantial advancements.
Early UAV systems were primarily designed for basic aerial imaging tasks and required extensive manual control.
In contrast, modern consumer-grade drones are equipped with advanced onboard sensors, reliable flight control algorithms, and stabilized camera systems.
The integration of gimbal stabilization, visual-inertial odometry, and autonomous navigation capabilities allows drones to acquire high-quality video data while maintaining stable and consistent viewpoints during flight.

These technological developments have significantly expanded the applicability of drones to vision-based tasks.
Contemporary UAV platforms support autonomous subject tracking, real-time video acquisition, and, in some cases, on-device processing.
Such capabilities make drones particularly well-suited for outdoor sports scenarios, where they can dynamically adapt their position and viewpoint to follow an athlete’s motion across large spatial areas, without requiring fixed infrastructure or intrusive setups.

At the same time, the use of drones introduces new challenges for human pose estimation.
Unlike static cameras, drones operate with continuously changing viewpoints, varying camera-to-subject distances, and non-negligible camera motion.
These factors result in significant scale variations, motion blur, and viewpoint-dependent appearance changes.
Furthermore, drone-based systems typically rely on a single moving camera, making 3D pose reconstruction inherently ill-posed due to the lack of multi-view geometric constraints.
Addressing these challenges is essential for enabling reliable human pose estimation from drone-acquired visual data.

\useImgWithCaption[]{Images/drone_tracking2.jpg}{0.5}{0.5}{An example of drone tracking}

\section{Drone-based Monocular 3D Human Pose Estimation}

The integration of drone technology with computer vision methods opens new opportunities for human pose estimation in outdoor sports scenarios.
Aerial viewpoints enable the acquisition of motion data in environments where fixed camera systems are impractical, providing large-scale and flexible coverage of complex athletic activities.
Compared to static sensing setups, drones can dynamically adapt their position to follow an athlete’s motion, capturing viewpoints that would otherwise be difficult or impossible to obtain.

In such unconstrained outdoor settings, traditional marker-based motion capture systems are not a viable solution.
The deployment of multiple calibrated cameras over large and potentially undefined areas is logistically challenging and economically costly.
Moreover, the use of physical markers is often incompatible with sports practice, as it may interfere with the athlete’s movement, alter natural performance, or be infeasible in competitive contexts.
For these reasons, markerless approaches are essential for practical sports analysis in real-world outdoor environments.

While multi-view configurations are commonly employed to resolve depth ambiguity in 3D pose estimation, their applicability in drone-based scenarios is limited.
Achieving a true multi-view setup would require the use of multiple coordinated drones, significantly increasing system complexity and cost.
Furthermore, maintaining diverse and complementary viewpoints becomes particularly challenging when athletes move at high speeds.
In such cases, multiple drones may be forced to follow similar trajectories, resulting in redundant views that provide limited additional geometric information.
The coordination and navigation of multiple aerial platforms also introduce non-trivial challenges that cannot be ignored in real-world deployments.

As a consequence, a single moving camera represents the most realistic and scalable configuration for drone-based sports analysis.
However, estimating 3D human pose from a single monocular view is a fundamentally ill-posed problem, as multiple three-dimensional poses can project to the same two-dimensional image.
Purely geometric approaches are therefore insufficient to recover reliable 3D structure from monocular observations alone.

To address this ambiguity, modern monocular 3D pose estimation methods rely on deep learning techniques to encode strong priors on human body shape, kinematics, and motion dynamics.
By learning these priors from large datasets, neural models can infer plausible three-dimensional poses that are consistent with the observed 2D evidence, even in the absence of explicit multi-view constraints.
Within this framework, drone-based monocular 3D human pose estimation emerges as a challenging yet practically relevant problem, balancing operational feasibility with the need for robust 3D motion understanding in outdoor sports contexts.

\useImgWithCaption[]{Images/drone_tracking.jpg}{0.5}{0.5}{An example of drone tracking}

\chapter{Domain background}

A keypoint is a semantically meaningful point in an image corresponding to a relevant part of an object or subject. In general, keypoints are characterized by the 
following properties:
\begin{enumerate}
    \item \textbf{Distinctiveness:} keypoints should be easily distinguishable from their surroundings, allowing them to be uniquely identified.
    \item \textbf{Invariance:} keypoints should be robust to common image transformations such as rotation, scaling, and changes in illumination.
    \item \textbf{Repeatability:} keypoints should be consistently detectable across different images or frames of the same object or scene.
\end{enumerate}
In the context of Human Pose Estimation, keypoints typically correspond to predefined human body joints or landmarks, such as the head, shoulders, elbows, wrists, hips, 
knees, and ankles.\\
Formally, a keypoint is represented by a spatial location, expressed as 2D coordinates \(x,y\) in the image plane or as 3D coordinates \(x,y,z\) in the physical space. 
In many approaches, a confidence score is also associated with each keypoint to indicate the reliability of the estimated location. The collection of keypoints 
belonging to a person defines a structured representation of the human body, commonly referred to as a pose or skeleton, which captures the geometric configuration of 
the subject and serves as a compact and interpretable representation for downstream tasks.

\useImgWithCaption[]{Images/keypoints.png}{1.0}{0.2}{Keypoints}



\section{2D Human Pose Estimation}
\label{sec:related_work_2d_hpe}
\subsection{Introduction}
2D Human Pose Estimation (HPE) has been a significant area of research in computer vision, with numerous approaches developed over the years. The 
task involves detecting and localizing human joints in images or videos, which is crucial for various applications such as activity recognition, 
human-computer interaction, and augmented reality. 2D keypoints extrapolation is often a preliminary step for 3D pose estimation, making it a 
foundational component in the broader field of human pose analysis.

The 2D pose estimation goals is to predict the coordinates of keypoints corresponding to human joints in a 2D image. With predict coordinates, we
mean the \((x, y)\) pixel locations of each joint in the image frame. The keypoints typically include joints such as the head, shoulders, elbows,
 wrists, hips, knees, and ankles.

When you are working with a 2D pose estimation model, the input is usually an image or a sequence of images (in the case of video). The output 
is a set of keypoints, containing the \((x, y)\) coordinates for each joint detected in the image, that represent the estimated pose of the person 
or people in the image.

\subsection{Evolution}
The evolution of 2D human pose estimation has seen significant advancements, particularly with the advent of deep learning techniques, I will
discuss only this type of approach in this section, that have revolutionized the field, with the advent of DeepPose.

The first tendency was to use Convolutional Neural Networks (CNNs) to directly regress the coordinates of keypoints from images. But this approach
faced challenges due to the complexity of human poses and variations in appearance, lighting, and occlusions. So researchers shifted towards 
heatmap-based methods, where the model predicts a probability distribution over possible locations for each keypoint, allowing for more robust 
localization that takes into account uncertainties in the predictions of joint positions.

TODO: DA FINIRE

\subsection{Principal approaches}
When we talk about 2D human pose estimation, we can divide the approaches into two main categories: single pose estimation and multi-pose
estimation.
In single pose estimation, the two main approaches are regression-based methods and heatmap-based methods. Regression-based methods directly 
predict the coordinates of keypoints from the input image, while heatmap-based methods generate heatmaps for each keypoint, indicating the 
likelihood of the keypoint's presence at each pixel location.
In multi-pose estimation, the two main approaches are top-down methods and bottom-up methods. Top-down methods first detect individual persons
in the image and then estimate the pose for each detected person separately. Bottom-up methods, on the other hand, first detect all keypoints 
in the image and then group them into individual poses. The choice between these approaches often depends on the specific application and the 
trade-offs between accuracy and computational efficiency, a top-down approach is generally more accurate but computationally expensive, while a
bottom-up approach is more efficient but may struggle with occlusions and overlapping poses.

In single pose video based estimation, is now common to use a single pose estimation model to extract 2D keypoints from each frame of the video,
and then thought about the temporal consistency between frames to improve the accuracy and stability of the pose estimation over time.
In multi pose video based estimation, there are two main strategies: online approaches and offline approaches. Online approaches process each 
frame sequentially, estimating poses in real-time as the video is being captured. Offline approaches, on the other hand, analyze the entire video 
after it has been recorded, allowing for more complex temporal modeling and potentially higher accuracy, as they can leverage information from 
future frames.

\subsection{Datasets, losses and metrics}




\subsection{Single Pose Estimation}


\subsection{Multi Pose Estimation}

\chapter{Problem Formulation}
The goal of this thesis is to estimate the 3D motion of athletes from monocular drone footage.\\ 
Specifically, the task consists in predicting the 3D HPE of an athlete over time, represented as a sequence of 3D joint coordinates.
Given a video acquired from a UAV-mounted camera, the objective is to recover the spatio-temporal configuration of the human body despite challenging acquisition conditions such as motion blur, occlusions, and dynamic camera motion.


\section{Input Representation}
The input to the model is an RGB video sequence :
\begin{equation}
    x \in \mathbb{R}^{T \times H \times W \times C}
\end{equation}
where $T$ is the temporal length of the video, $H$ and $W$ are the spatial dimensions of each frame and $C$ is the number of channels.\\
The video is assumed to be captured by a moving drone-mounted monocular camera, resulting in non-static backgrounds and varying viewpoints.


\section{Output Representation}
The output is a sequence of 3D human poses, one for each video frame. 
Each pose is represented as a set of $J$ 3D joint coordinates:
\begin{equation}
    y \in \mathbb{R}^{T \times J \times P}
\end{equation}
where $T$ is the temporal length of the sequence and $P = 3$ represents for each joint its $(x, y, z)$ coordinates in a root-relative reference frame, obtained by translating all joints such that the root joint is located at 
the origin.


\section{Task Definition}
The problem addressed in this thesis is to learn a function $f_\theta$, parameterized by $\theta$, that maps the input video to a sequence of 3D human poses:
\begin{equation}
    f_\theta: x \in \mathbb{R}^{T \times H \times W \times C} \rightarrow y \in \mathbb{R}^{T \times J \times P}
\end{equation}


\section{Learning Objective}
Given a training set of video sequences and corresponding ground-truth 3D poses, the learning objective is to minimize the discrepancy between predicted and ground-truth joint positions.
This is typically achieved by minimizing the \textit{MPJPE} loss (see \ref{mpjpe}) over all joints and time steps:
\begin{equation}
    \text{MPJPE} = \frac{1}{T} \frac{1}{J} \sum_{t=1}^{T} \sum_{j=1}^{J} \| \hat{\mathbf{y}_{t,j}} - \mathbf{y}_{t,j} \|
\end{equation}


\section{Challenges and Constraints}
Estimating 3D human pose from drone footage presents several challenges:
\begin{itemize}
    \item \textbf{Motion blur and low resolution}, caused by fast camera and subject motion;
    \item \textbf{Self-occlusions and partial visibility} of the human body;
    \item \textbf{Dynamic camera motion}, which prevents the use of static background assumptions;
    \item \textbf{Limited availability of labeled datasets}, as only a few public UAV-based 3D human pose datasets exist.
\end{itemize}
These factors make the problem significantly more challenging than standard 3D HPE from fixed ground cameras.



\chapter{Proposed Method}
Now I will describe the proposed method in detail. The proposed method consists of several key components that work together to achieve the desired outcome. I proposed the network \textbf{Brancher} that is a two-stage network. Given a
raw-video, the model detects the human in the scene, extracts his keypoints and then lifts them to 3D. It is proposed a self-supervised approach in a way that after a first training of the model on a labeled dataset, the model 
can continuos learning with unlabeled videos, in a way to fill the gap with the scarcity of 3d labeled data captured in the wild by drone. The model uses a spatio-temporal attention module to better capture the temporal dependencies 
between frames and the spatial relationships between keypoints. Then different branch are proposed to better capture the different aspects of the human motion.

First a simple baseline model is presented.

\section{Preprocess}
When dealing with RGB videos, it is fundamental to preprocess the input data, since otherwise the model will explode in ram and also the training time will be really high and long.
To do so first the video are need to be resized to a fixed resolution in a way to have a fixed input size and to not explode the ram memory with full hd videos.
Then a human detector is used to detect the human in the scene and crop around him. In this way the model will focus only on the human and not on the background that is not important for the task. The detector used is the YoloV12
that is a fast and accurate object detector.
Then a 2D keypoint detector pretrained is used to extract the 2D keypoints from the cropped human frames. It extracts a heatmap for each keypoint and then the argmax is used to extract the 2D coordinates and the confidence score of
each keypoint. The keypoint detector used is ViTPose.
Finally the 2D keypoints are normalized in range [-1,1] with respect to the frame size in a way to avoid the exploding gradients problem during the training.

The preprocess video began a sequence of 2D keypoints that will be the input of the model.

\section{Baseline}
A baseline model is proposed to have a starting point for the development of the final model and for the ablation study to compare the different components of the final model.
The baseline model takes as input a sequence of 2D keypoints preprocessed as described in the previous section, and outputs the 3D keypoints of the human in the scene in root-relative coordinates.
The baseline model embedded the input 2D keypoints with a linear layer to increase the feature dimension, and then they are encoded with a standard positional encoding to add the temporal information, with a classic sinusoidal encoding.
Then the encoded features are passed through a series of Transformer Encoder layers to capture the temporal dependencies between frames. In the baseline the attention mechanism is implemented only between frames, without considering 
the spatial relationships between keypoints.
Finally a MLP head is used to regress the 3D keypoints from the encoded features. A tanh activation function is applied to capture the non-linearities.
The output 3D keypoints are in root-relative coordinates, meaning that the root joint (usually the pelvis) is at the origin (0,0,0).

The baseline model is trained with a supervised loss function, using the Mean Per Joint Position Error (MPJPE) between the predicted 3D keypoints and the ground truth 3D keypoints.


\section{Proposed Network Architecture}
The proposed network architecture is an improvement over the baseline model, with the addition of a spatial-temporal attention module and a multi-branch architecture to better capture the different aspects of the human motion.
In other is self-supervised way the model is able to learn also from unlabeled videos, in a way to fill the gap with the scarcity of 3d labeled data captured in the wild by drone. 

\subsection{Spatial Temporal Attention Module}
The spatial-temporal attention module is designed to capture both the temporal dependencies between frames and the spatial relationships between keypoints.
It is inspired by the work of MotionAGFormer.

\subsection{Multi-Branch}
The multi-branch architecture, inspired by SelfPose, consists of different branches that are specialized in capturing different aspects of the human motion, in a way to improve the performance of a model that has to deal with complex 
human motion taken from oblique angles by drone and in the wild.
Each branch is composed by an MLP head that takes as input the encoded features from the spatial-temporal attention module and outputs a specific aspect of the human motion using different loss functions to better capture the 
different aspects of the motion.

\paragraph{Pose Branch}
This branch is responsible for capturing the overall pose of the human in the scene. It outputs the 3D keypoints of the human in root-relative coordinates. It is an MLP head that takes as input the encoded features from the 
spatial-temporal attention module and outputs the 3D keypoints. A tanh activation function is applied to capture the non-linearities. The loss function used for this branch are the supervised MPJPE and a self-supervised 
bone length loss that ensures the consistency of the bone lengths penalizing strong variations in the length of the bones between consecutive frames, a velocity loss that ensures the smoothness of the motion penalizing strong 
variations in the velocity between consecutive frames, and an acceleration loss that ensures the smoothness of the motion penalizing strong variations in the acceleration between consecutive frames.

\begin{equation}
    \mathcal{L}_{bone} = \frac{1}{B} [\text{Var}_{t}(l_{bone})]
\end{equation}
\begin{equation}
    \mathcal{L}_{vel} = \frac{1}{(T-2)J} \sum_{t=2}^{T-1} \sum_{j=1}^{J} \| (\hat{\mathbf{y}_{t+1,j}} - \hat{\mathbf{y}_{t,j}}) - (\hat{\mathbf{y}_{t,j}} - \hat{\mathbf{y}_{t-1,j}}) \|_2
\end{equation}
\begin{equation}
    \mathcal{L}_{acc} = \frac{1}{(T-3)J} \sum_{t=2}^{T-2} \sum_{j=1}^{J} \| (\hat{\mathbf{y}_{t+1,j}} - 2\hat{\mathbf{y}_{t,j}} + \hat{\mathbf{y}_{t-1,j}}) - (\hat{\mathbf{y}_{t,j}} - 2\hat{\mathbf{y}_{t-1,j}} + \hat{\mathbf{y}_{t-2,j}}) \|_2
\end{equation}


\paragraph{Rotation Branch}
This branch is responsible to compute the rotations of each joint in order to have a more accurate representation of the human motion. It used a MLP head that takes as input the encoded features from the spatial-temporal attention module 
and uses the 6D rotation representation to output the rotations of each joint. The Gram-Schmidt orthogonalization is used to convert the 6D representation to a valid rotation matrix. This branch taken in input also the output of the 
Pose Branch and calculate from these points the local joint rotations with respect to the T-Pose. So this is a self-supervised branch that does not need 3D ground truth data to be trained. The loss function used for this branch is 
the geodesic loss that measures the distance between two rotations.

\begin{equation}
    \mathcal{L}_{geo} = \frac{1}{B} \text{arccos} \left( \frac{\text{tr}(R_{\text{pred}} R_{\text{target}}^T) - 1}{2} \right)
\end{equation}

\paragraph{Uncertainty Branch}
This branch is responsible for estimating the uncertainty of the predicted 3D keypoints. It used a MLP head that takes as input the encoded features from the spatial-temporal attention module and reconstructs the heatmaps of each
joints.


The different loss are weighted with learnable parameters during the training phase using the homoscedastic uncertainty method.
\begin{equation}
    \mathcal{L}_{total} = \sum_{i} \frac{1}{\sigma_i^2} \mathcal{L}_i + \log \sigma_i
\end{equation}


\chapter{Implementation Details}
This section describes the practical implementation of the proposed framework. I first detail the baseline model, including its architecture, training procedure, and 
inference pipeline. All implementation choices are reported to ensure reproducibility of the experiments. All code is implemented in Python using the PyTorch library.
All the hyperparameters are selected based on validation performance using the library Optuna and kept constant across all experiments unless otherwise specified. 
The experiments were conducted on a machine with TODO

\section{Baseline Implementation}
The \hyperref[sec:baseline]{baseline model} serves as the reference architecture upon which the proposed method is built. It takes as input a temporal sequence of 2D human poses $[B, T, J, 3]$,
where $B$ is the batch size equal to 32, $T$ is the number of frames in the sequence equal, and $J$ is the number of joints equal to 17, and 3 represents the $(x, y, confidence)$ coordinates of each 
joint and predicts the corresponding 3D pose representation, giving as output a tensor of shape $[B, T, J, 3]$, where the last dimension represents the $(x, y, z)$ 
coordinates of each joint in 3D space.

\subsection{Architecture}
The baseline model is based on a Transformer-based architecture designed to model temporal dependencies across pose sequences. Each input pose is first embedded into a 
latent feature space of dimension 512. The resulting embeddings are then processed by a stack of Transformer encoder layers.

Specifically, the Transformer module consists of 3 encoder layers, each employing multi-head self-attention with 8 attention heads. This design allows the model to 
capture both short-term and long-range temporal correlations within the input sequence. The output of the Transformer is finally mapped to the target pose 
representation through a linear projection layer.

No architectural modifications are introduced with respect to the baseline configuration, which is kept fixed throughout all experiments to provide a fair comparison 
with the proposed approach.

\begin{table}[t]
    \centering
    \caption{Baseline model architecture. The network takes as input a temporal sequence of 2D poses and outputs the corresponding pose representation.}
    \label{tab:baseline_architecture}
    \renewcommand{\arraystretch}{1.2}
    \begin{tabular}{lccc}
    \hline
    \textbf{Layer} & \textbf{Input Shape} & \textbf{Output Shape} & \textbf{Params} \\
                   & $(32, 243, 17, 3)$ & $(32, 243, 17, 3)$ &  \\
    \hline
    Linear (Input Embedding) & $(32, 243, 51)$ & $(32, 243, 512)$ & 26,624 \\
    \hline
    Positional Encoding & $(32, 243, 512)$ & $(32, 243, 512)$ & -- \\
    Transformer Encoder (3 layers, 8 heads) & $(32, 243, 512)$ & $(32, 243, 512)$ & 9,457,152 \\
    \hline
    Linear & $(32, 243, 512)$ & $(32, 243, 512)$ & 262,656 \\
    Tanh & $(32, 243, 512)$ & $(32, 243, 512)$ & -- \\
    Linear (Pose Head) & $(32, 243, 512)$ & $(32, 243, 51)$ & 26,163 \\
    \hline
    \textbf{Total} &  &  & \textbf{9,772,595} \\
    \hline
    \end{tabular}
\end{table}


\subsection{Training Procedure}
The baseline model is trained using the MPJPE as the optimization objective. Training is initially performed on the SportsPose dataset, which provides labeled pose 
sequences suitable for learning temporal motion patterns. Subsequently, the model is fine-tuned on the AirPose dataset to adapt it to the target domain.

The training process is carried out for a maximum of 100 epochs. The Adam optimizer is employed with an initial learning rate of $5 \times 10^{-4}$ and a weight decay of 
0.01. The learning rate is scheduled using a exponential decay with a factor of 0.99 applied every epoch.
A warm-up phase of 2 epochs is applied at the beginning of training to stabilize optimization. Early stopping is used with a patience of 9 epochs based on 
validation performance.

\subsection{Inference Pipeline}
During inference, the input to the system is a RGB video sequence represented as a tensor of shape $[T, H, W, 3]$, where 
For each frame, a human detection stage is first performed using YOLO to localize the subject of interest.

The detected human bounding box is cropped and resized to a fixed resolution of $256 \times 192$ pixels. The resulting image is then passed to a ViTPose model to 
perform 2D pose estimation, producing a set of 2D keypoints for each frame.

The extracted 2D keypoints are normalized to the range $[-1,1]$ using the original frame dimensions.
These normalized 2D poses are finally provided as input to the baseline network, which predicts the corresponding output pose sequence.


\chapter{Experiments and Results}

\section{Evaluation Setting}

Initially, I intended to evaluate the proposed method on the UAV dataset introduced by AirPose \cite{saini2022airpose}, as the authors provide calibration matrices for the two drones used during acquisition. In principle, this makes 
it possible to reconstruct 3D joint coordinates through triangulation and use them as pseudo ground truth for quantitative evaluation.

However, in practice, the triangulated 3D poses turned out to be excessively noisy. The reconstructed joints were highly unstable due to synchronization inaccuracies, imperfect detections, and sensitivity of triangulation to small 
2D localization errors. As a consequence, the resulting 3D annotations were not reliable enough to produce meaningful quantitative comparisons. The evaluation metrics computed on these reconstructions were inconsistent and did not 
reflect the true quality of the predicted poses.

For this reason, I performed the quantitative evaluation on the test split of AthletePose3D \cite{yeung2025athletepose3d}, where accurate 3D ground-truth annotations are available. This allows a reliable and reproducible comparison 
between models under controlled conditions.

\section{Evaluation Protocol}

During evaluation, 2D keypoints are first extracted and normalized following the preprocessing procedure described in the Section~\ref{sec:preprocess}. Temporal windows of 81 consecutive frames are fed to the network using a 
sliding-window strategy consistent with training.

Predicted 3D poses are root-centered before metric computation. When enabled, test-time augmentation (TTA) through horizontal flipping is applied, and final predictions are obtained by averaging the outputs of the original and 
flipped inputs.

\section{Evaluation Metrics}

To assess performance, I report Mean Per Joint Position Error (MPJPE) and Procrustes-aligned MPJPE (P-MPJPE), both measured in millimeters.

MPJPE computes the average Euclidean distance between predicted and ground-truth 3D joint coordinates. P-MPJPE applies a rigid Procrustes alignment before error computation, removing global rotation and translation differences. 
This metric isolates structural pose accuracy independently of global orientation.


\section{Quantitative Results}

\begin{table}[H]
\centering
\caption{Quantitative comparison performed on the AthletePose3D test set.}
\label{tab:results}
\begin{tabular}{lcc}
\hline
Model & MPJPE (mm) & P-MPJPE (mm) \\
\hline
MotionBERT & -- & -- \\
MotionAGFormer & 160.6560 & 86.1506 \\
MotionAGFormer Finetuned & 14.1359 & 11.0060 \\
Baseline & 45.2764 & 4.3332 \\
Brancher & 30.9958 & 27.8366 \\
Brancher Self Supervised Finetuned & -- & -- \\
\hline
\end{tabular}
\end{table}

Table~\ref{tab:results} reports the quantitative performance of all evaluated models on the AthletePose3D test split.

The original MotionAGFormer model \cite{mehraban2024motionagformer}, trained on Human3.6M datasets \cite{zhu2023motionbert}, exhibits the highest error, with an MPJPE of 160.66 mm and a P-MPJPE of 86.15 mm. After fine-tuning on 
AthletePose3D, MotionAGFormer Finetuned achieves a substantial reduction in error, reaching 14.14 mm MPJPE and 11.01 mm P-MPJPE.

The Baseline model obtains an MPJPE of 45.28 mm and a P-MPJPE of 4.33 mm. Compared to MotionAGFormer Finetuned, it presents a higher MPJPE but a lower Procrustes-aligned error.

The proposed Brancher architecture achieves an MPJPE of 31.00 mm and a P-MPJPE of 27.84 mm, improving over the Baseline in terms of MPJPE while yielding a higher P-MPJPE.

TODO: da finire e dire che questa era una prima analisi dei risultati basata unicamennte sui numeri.


\section{Qualitative Analysis}


\chapter{Conclusions}
The objective of this thesis was to develop a deep learning model capable of performing 3D Human Pose Estimation from sports videos captured by a drone. To address this problem, I proposed Brancher, a 
single-view architecture designed to operate under aerial viewpoints and challenging motion dynamics. The model was conceived to satisfy the requirements defined in the problem formulation while 
contributing to the current research landscape, where single-view 3D HPE methods specifically tailored to UAV footage remain limited.

Brancher integrates spatio-temporal feature manipulation with uncertainty-aware learning. By jointly modeling spatial structure and temporal dynamics, and by optimizing 2D consistency, rotational alignment, 
and prediction uncertainty, the model mitigates the effects of occlusions, viewpoint variations, and noisy pseudo-labels. This design enables the reconstruction of complex athletic 3D motions from 
challenging drone footage while maintaining structural coherence and temporal stability.

Both quantitative and qualitative evaluations demonstrate that Brancher achieves solid performance compared to existing baselines. The architecture benefits from a state-of-the-art backbone, yet extends it 
through task-specific modifications that improve motion plausibility and robustness under aerial observation. Importantly, the model remains computationally lightweight, with approximately 5 million 
parameters. This relatively compact design opens the possibility for future onboard deployment directly on UAV platforms, enabling real-time 3D motion analysis in sports scenarios.

Although the model is subject to the limitations discussed in the previous chapter, particularly those related to the absence of ground-truth 3D annotations and the intrinsic challenges of purely 
self-supervised adaptation, it represents a meaningful step toward practical 3D human motion reconstruction from aerial video. I developed a system capable of balancing multiple optimization objectives to 
address a complex and underexplored problem domain, namely 3D HPE under UAV-based acquisition. Beyond its methodological contribution, Brancher is intended as a step toward practical motion analysis in 
sports performance, with the long-term goal of supporting performance optimization and injury prevention through more accessible 3D biomechanical reconstruction.

This research project evolved progressively, starting from keypoint detection and 2D HPE, extending to 3D HPE, and ultimately reaching the more demanding setting of 3D reconstruction from 
UAV footage. The path toward Brancher required navigating a vast body of literature, analyzing hundreds of publications to identify both methodological gaps and promising directions. 

The development process involved extensive experimentation, iterative model redesign, and long training cycles performed on high-performance GPU servers. Multiple architectural configurations were tested 
before converging to the final formulation. This journey was not only a technical endeavor but also an exercise in critical reasoning, persistence, and methodological rigor.

Working on this problem highlighted how research is rarely a linear process. Instead, it requires continuous refinement of ideas, empirical validation, and acceptance of both progress and limitations. The 
transition from controlled indoor benchmarks to unconstrained aerial sports footage exemplifies the gap between theoretical performance and real-world applicability. Bridging this gap has been one of the 
central motivations of this work.

\section{Future Work}
A first and fundamental direction for future research is the creation of a dedicated UAV-based 3D Human Pose Estimation dataset. Such a dataset could be developed either through high-fidelity simulation environments or through 
real-world acquisition using multiple synchronized drones to reconstruct 3D pose with high precision. \\
The performance of deep learning models is strongly dependent on the quality and representativeness of the training data. In the context of aerial sports analysis, the lack of 3D-annotated UAV footage remains one of the main 
bottlenecks. The availability of a large-scale dataset specifically tailored to sports scenarios captured from drones would significantly enhance the training process and provide stronger supervision for models such as Brancher. 
This would also reduce the reliance on purely self-supervised adaptation strategies and mitigate the geometric ambiguities inherent to monocular reconstruction.

A second promising direction concerns the development of models capable of learning from datasets annotated with heterogeneous joint conventions. In practice, many 2D and 3D datasets adopt different skeletal definitions 
(e.g., COCO, Human3.6M, and others), which limits the possibility of jointly exploiting multiple data sources.\\ 
While Human3.6M \cite{ionescu2013human3} has become the de facto standard for 3D supervision due to its size and popularity, sports-related datasets are often smaller and follow alternative annotation schemes. Designing architectures that can align, map, 
or dynamically adapt across different skeletal representations would allow training on a broader range of datasets. In data-scarce domains such as UAV sports footage, this capability could substantially improve generalization and robustness.
\useImgWithCaption[]{Images/different_notations.png}{0.25}{0.25}{This picture taken from the work of Sarandi et al. \cite{sarandi2023learning}, shows the diversity of joint definitions across different datasets, in a work that addresses the problem of learning from heterogeneous annotations.}

Finally, in the current era of generative artificial intelligence, diffusion-based models offer an intriguing opportunity for advancing single-view 3D HPE. A possible research direction would involve 
training a diffusion model that takes normalized 2D joint coordinates as input and generates additional plausible viewpoints of the same pose.\\
By leveraging geometric constraints such as epipolar consistency or learned cross-view priors, synthetic multi-view representations could be constructed from monocular input. This would effectively 
transform a single-view reconstruction problem into a richer multi-view setting, where triangulation or multi-view consistency losses could be applied.\\ 
Such an approach could bridge the gap between monocular ambiguity and multi-view geometric constraints, opening new possibilities for 3D human motion reconstruction from UAV footage without requiring 
physically deployed multi-drone systems.
\useImgWithCaption[]{Images/diffusion_model.png}{1.0}{0.28}{This image, taken from the work of Li et al. \cite{li2025lifting}, shows a schematic representation of a diffusion model for generating multi-view 3D human poses from monocular 2D joint coordinates.}
\newpage
\null\vspace{1cm}
With this final section, my thesis work comes to an end, marking also the conclusion of my academic journey at Politecnico di Milano. This work represents the culmination of years of study, research, and personal growth.

Thank you for taking the time to read my Master's thesis.

\vspace{0.5cm}

\noindent
Sincerely,\\
Angelo


%

%-------------------------------------------------------------------------
%	BIBLIOGRAPHY
%-------------------------------------------------------------------------

\addtocontents{toc}{\vspace{2em}} % Add a gap in the Contents, for aesthetics
\bibliography{Thesis_bibliography} % The references information are stored in the file named "Thesis_bibliography.bib"

%-------------------------------------------------------------------------
%	APPENDICES
%-------------------------------------------------------------------------

\cleardoublepage
\addtocontents{toc}{\vspace{2em}} % Add a gap in the Contents, for aesthetics
\appendix
\chapter{Appendix A}
\section{Drone-Captured Self-Supervised Dataset}
\label{appendix:drone_dataset}

To support the self-supervised training stage of Brancher, a dedicated drone-based dataset was collected using a commercial \textit{Parrot Bebop 1} quadcopter.
\useImgWithCaption[]{Images/bebop.jpg}{0.25}{0.15}{The Parrot Bebop 1 drone used for data collection.} 
The goal of this acquisition was to obtain monocular aerial videos of human motion under realistic drone viewpoints, without relying on ground-truth 3D annotations.

The dataset consists of \textbf{698 video sequences}, each with a fixed temporal length of \textbf{81 frames}. All videos were recorded at a spatial resolution of $\mathbf{H=352}$ and $\mathbf{W=640}$ 
pixels. This fixed-length design ensures full compatibility with the temporal window used during training.

Data were collected across \textbf{five different locations} and involve \textbf{five distinct subjects}. The locations include varied outdoor environments, allowing changes in background, 
illumination, and scene geometry. The subjects were instructed to perform movements that generate highly diverse body configurations, covering a broad range of poses and limb articulations. This 
variability is crucial to avoid pose bias and to improve generalization under aerial viewpoints.

Importantly, each video focuses on a \textit{single person} captured from a drone perspective, resulting in monocular, single-view pose sequences. No 3D ground-truth annotations were recorded. The 
dataset was explicitly designed for \textbf{self-supervised training}, enabling the model to adapt to drone-specific viewpoints and motion dynamics without relying on supervised 3D labels.

This drone-captured dataset plays a central role in the second training stage of Brancher, where the network is optimized in a fully self-supervised manner to better generalize to real-world UAV 
scenarios.

\useImgWithCaption[]{Images/drone_dataset.jpg}{1.0}{0.6}{Some samples from the drone-captured dataset, showcasing the diversity of subjects, poses, and environments.}

\chapter{Appendix B}
\label{app:qualitative}
In this appendix, I present additional qualitative results that complement the quantitative findings discussed in the main body of the thesis. These results provide further insights into the behavior of the model under study, and 
help to illustrate key concepts and phenomena that were observed during my experiments.
\useImgWithCaption[]{Images/appendice_cortile.png}{1.0}{0.4}{Some examples where the subject interacts with some objects.}
\useImgWithCaption[]{Images/appendice_cortile_c.png}{1.0}{0.8}{Some examples from the custom dataset.}
\useImgWithCaption[]{Images/appendice_cortile_i.png}{1.0}{0.8}{Some examples from the custom dataset.}
\useImgWithCaption[]{Images/appendice_parco.png}{1.0}{0.8}{Some examples from the custom dataset.}
\useImgWithCaption[]{Images/appendice_iceskating.png}{1.0}{0.8}{Some examples of ice skating.}
\useImgWithCaption[]{Images/appendice_sci.png}{0.65}{0.55}{Some examples of sky.}

% LIST OF FIGURES
\listoffigures

% LIST OF TABLES
\listoftables

% LIST OF SYMBOLS
% Write out the List of Symbols in this page
%\chapter*{List of Symbols} % You have to include a chapter for your list of symbols


% ACKNOWLEDGEMENTS
\chapter*{Acknowledgements}
Arrivato alla fine di questo lungo viaggio iniziato nel 2019, non posso che ringraziare tutte le persone che mi hanno accompagnato in questo percorso, sia dal punto di vista accademico che personale e che mi hanno aiutato a 
raggiungere questo traguardo e superare gli ostacoli che si sono presentati lungo il cammino.

In primo luogo, vorrei ringraziare i miei genitori, che mi hanno dato la possibilità di studiare e mi hanno sempre sostenuto in ogni mia scelta. Grazie per avermi insegnato il valore del duro lavoro e per avermi sempre incoraggiato a 
perseguire i miei sogni.

Un grazie alla mia ragazza, che mi ha supportato in ogni momento, anche nei periodi più difficili. Grazie per la tua pazienza, il tuo amore e la tua comprensione.

Un ringraziamento speciale va al mio relatore e al mio co-relatore, per la vostra guida, il vostro supporto e i vostri preziosi consigli durante tutto il percorso di ricerca. Grazie per aver creduto in me e per il supporto dato a 
questo progetto. Grazie anche a tutti i membri del team di ricerca, per la collaborazione e il supporto durante il lavoro di tesi.

Grazie anche a tutti i miei amici e ai miei colleghi, per il supporto, le risate e i momenti di svago che mi hanno aiutato a mantenere un equilibrio tra studio e vita personale.

Infine, vorrei ringraziare tutti coloro che hanno contribuito in qualche modo a questo progetto e alle riprese del dataset UAV, anche se non menzionati direttamente. Grazie a tutti voi per aver reso possibile questo traguardo.

È stato un viaggio impegnativo ma gratificante, e sono grato a tutti coloro che mi hanno accompagnato lungo il cammino. Grazie di cuore a tutti voi!

Un ringraziamento finale va invece a me stesso, per la perseveranza, la determinazione e la passione che mi hanno spinto a superare le difficoltà. Dalle prime lezioni di Analisi 1 fino al testing del codice finale di questa tesi,
è stato un percorso di crescita personale e accademica che mi ha permesso di raggiungere questo importante traguardo. Grazie a me stesso per aver creduto in me e per aver lavorato duramente per arrivare fin qui.

\cleardoublepage

\end{document}
